
\chapter{Présentation du CNRS et de l’IN2P3}
\label{chap:cnrs-in2p3}
\justifying
\sloppy

Le \textbf{Centre National de la Recherche Scientifique (CNRS)} occupe une place centrale dans le paysage de la recherche publique française. Organisme pluridisciplinaire, il a pour mission de produire des connaissances, de mobiliser les différentes disciplines scientifiques et d’éclairer les grands défis contemporains, en lien avec les autres acteurs de la recherche \cite{cnrs2025lecnrs}. Cette organisation permet d’articuler recherche fondamentale, ingénierie et coopération avec les universités, les laboratoires et les grandes infrastructures scientifiques.

Dans les domaines relevant de la physique et de l’instrumentation, le CNRS joue un rôle déterminant non seulement par ses contributions scientifiques, mais aussi par sa capacité à fédérer des moyens humains et techniques autour de projets de long terme. Cette dimension est particulièrement visible dans les activités liées à la conception de détecteurs, à l’électronique associée, au traitement du signal et au développement de circuits intégrés spécialisés pour des environnements expérimentaux exigeants.

\section{L’IN2P3 : un institut tourné vers la physique subatomique et l’instrumentation}

Parmi les instituts du CNRS, l’\textbf{IN2P3 (Institut National de Physique Nucléaire et de Physique des Particules)} joue un rôle majeur dans la coordination de la recherche française en physique subatomique. Son champ d’action couvre notamment la physique des particules, la physique nucléaire, l’astrophysique et la cosmologie, ainsi que le développement des moyens expérimentaux nécessaires à ces disciplines. À ce titre, l’IN2P3 ne se limite pas à une activité purement théorique : il contribue également à la conception d’instrumentations complexes, de chaînes d’acquisition et de systèmes électroniques adaptés à des besoins de mesure de très haute précision \cite{in2p32025presentation}.

L’activité de l’IN2P3 s’appuie sur un réseau de laboratoires, de plateformes technologiques et de collaborations nationales et internationales. Cette structuration favorise la conduite de projets ambitieux, dans lesquels les performances des détecteurs dépendent directement de la qualité de l’électronique développée en amont. Les problématiques de chronométrie fine, de lecture rapide, de robustesse et d’intégration sur \textit{ASIC} occupent ainsi une place importante dans plusieurs thématiques de recherche portées par l’institut \cite{in2p32025presentation}.

L’alternance présentée ici s’inscrit dans cet environnement scientifique et technologique. Réalisée au sein de l’\textbf{IP2I}, elle prend place dans un contexte où l’électronique de détection constitue une brique essentielle de la chaîne instrumentale. Le travail mené autour du \textbf{TDC} destiné au projet \textbf{SPADMIC} relève précisément de cette logique : traduire un besoin de mesure temporelle fine en une architecture intégrable sur circuit, compatible avec des contraintes élevées de précision, de lecture et de fiabilité.

\clearpage
